\chapter{Por que o Java é assim}
\label{cap:por-que}

\proximasaulas

Por que o Java e a JVM são assim, e o que isso muda no código que você
escreve? Este capítulo responde a essa pergunta com foco em concorrência:
muitos \termo{requests}{Requisições que chegam ao servidor, cada uma
exigindo trabalho do seu programa.}, memória compartilhada entre
\termo{threads}{Linhas de execução dentro do mesmo processo, que
compartilham a memória.} e vários núcleos.
Cada seção segue a mesma forma: o problema da época, a decisão tomada, o
que mudou depois, o que isso muda no seu código e onde aprofundar.
As listagens de bytecode são saída real do \texttt{javap} do JDK~25
sobre os exemplos de \texttt{examples/}. Se o \texttt{javap} for
novidade, o capítulo~\ref{cap:lendo-bytecode} ensina a ler essas
listagens; volte aqui depois.

Afirmações históricas são conferidas nas segundas edições da JLS
(\citetitle{livro:gosling2000-jls2}) e da JVMS
(\citetitle{livro:lindholm1999}). Afirmações sobre o estado atual da
linguagem e da máquina virtual são conferidas na JLS e na JVMS do Java
SE~25, \citetitle{doc:jls25} e \citetitle{doc:jvms25}.

\section{Os objetivos de 1995}
\label{sec:objetivos-1995}

O problema da época era distribuir código por redes de máquinas
diferentes. O white paper \citetitle{paper:java-language-environment}
(maio de 1996, de Gosling e McGilton, lido hoje via Wayback) afirma que a
tecnologia deve permitir \enquote{secure, high performance, and highly
robust applications on multiple platforms in heterogeneous, distributed
networks}. O mesmo texto resume o sistema como simples, familiar,
orientado a objetos, \emph{multithreaded} e interpretado, para máxima
portabilidade e capacidades dinâmicas.

\paragraph{A decisão.}
As metas de segurança, robustez e neutralidade de arquitetura (§1.2.2 e
§1.2.3 do white paper) levaram a quatro escolhas que o resto do capítulo
retoma. Primeiro, o compilador gera \termo{bytecode}{Conjunto de
instruções da JVM, independente de processador. O \texttt{javac} o
grava em arquivos \texttt{.class}.}: um formato intermediário neutro,
pensado para transportar código a várias plataformas (§1.2.3). Segundo,
o bytecode foi desenhado para ser fácil de interpretar em qualquer
máquina ou de traduzir dinamicamente para código nativo quando o
desempenho pedir (§4.1.1). Terceiro, o código que chega pela rede não é
confiável: o sistema de execução o submete à verificação de bytecode
(§6.3), tratada na seção~\ref{sec:verificacao}. Quarto, o modelo de
memória não tem tipos ponteiro definidos pelo programador, nem aritmética
de ponteiros, e tem coleta automática de lixo (§1.2.2); o white paper
também diz que programadores Java não conseguem forjar ponteiros para a
memória (§6.1). Para a coleta, o texto descreve um \termo{garbage
collector}{Coletor de lixo: componente da JVM que libera os objetos que
o programa não alcança mais.} em segundo plano, numa thread de baixa
prioridade (§2.1.7). Já as threads são parte da linguagem, preemptivas e,
conforme a plataforma, divididas por fatias de tempo (§7.1), assunto da
seção~\ref{sec:monitores}.

A neutralidade de arquitetura tem um custo declarado: a JVM é descrita
como máquina abstrata baseada principalmente na interface POSIX
(§1.2.3), o que pressupõe suporte a multithreading na plataforma de
destino.

\paragraph{O que mudou depois.}
O formato do bytecode continuou o mesmo em espírito, mas a forma de
verificá-lo mudou (seção~\ref{sec:verificacao}) e o carregamento de
classes ficou mais elaborado (seção~\ref{sec:class-loading}). O white
paper chama o carregador de bytecode de próxima linha de defesa depois
do verificador (§6.2).

\paragraph{O que isso muda no seu código.}
Vale para qualquer JDK, inclusive o 25: você compila uma vez e o mesmo
\texttt{.class} roda onde houver JVM. Você não faz aritmética de
ponteiros nem libera objetos manualmente; em troca, o tempo de coleta e
o comportamento de threads entram no seu raciocínio sobre desempenho,
temas das seções seguintes.

\paragraph{Onde aprofundar.}
O formato do bytecode é lido em detalhe no capítulo~\ref{cap:lendo-bytecode},
e o carregamento, a ligação e a inicialização estão na
seção~\ref{sec:class-loading}.

\section{A verificação do bytecode}
\label{sec:verificacao}

\section{Monitores e \texttt{synchronized}}
\label{sec:monitores}

\paragraph{O problema.}
Na JLS 2ª ed., as threads executam código que opera sobre valores e
objetos de uma memória principal compartilhada. A linguagem apoia a escrita de
programas que, mesmo concorrentes, tenham comportamento determinístico,
e para isso fornece mecanismos de sincronização (cap.~17, abertura). O
mesmo texto descreve a alternativa a evitar: \enquote{spinning}, travar
e destravar um objeto repetidamente para ver se o estado interno mudou,
o que consome esforço computacional.

\paragraph{A decisão.}
A resposta da JLS 2ª ed. é o \termo{monitor}{Mecanismo que deixa uma
única thread por vez executar a região de código que ele protege.}:
\enquote{a high-level mechanism for allowing only one thread at a time
to execute a region of code protected by the monitor}. O comportamento
do monitor é explicado em termos de \termo{locks}{Travas. Uma thread
adquire o lock antes de entrar na região protegida e o libera ao
sair.}, um por objeto (cap.~17, abertura). A linguagem não oferece
ações separadas de travar e destravar: o bloco \texttt{synchronized} e
o método \texttt{synchronized} as executam aos pares. Um método de
instância trava o lock de \texttt{this}, e um método \texttt{static}
trava o do objeto \texttt{Class} da classe (§17.13). Os métodos
\texttt{wait}, \texttt{notify} e \texttt{notifyAll} de \texttt{Object}
dão suporte à transferência eficiente de controle entre threads: em vez
de girar, uma thread pode se suspender com \texttt{wait} até outra
acordá-la com \texttt{notify}
(cap.~17, abertura).

Cada objeto tem, além do lock, um \termo{wait set}{Conjunto das
threads que chamaram \texttt{wait} no objeto e ainda não foram
removidas dele.} (§17.14). Segundo a §17.14, \texttt{wait} só deve ser
chamado por quem já detém o lock do objeto. A chamada põe a thread no
wait set e executa tantas liberações quantos forem os travamentos
pendentes. A thread fica dormente até acontecer uma de três coisas: um
\texttt{notify} que a escolha arbitrariamente, um \texttt{notifyAll} ou
o fim do tempo limite, se houver. Depois ela trava o objeto de novo e
volta de \texttt{wait} com o lock no mesmo estado em que estava.

A JVMS 2ª ed. traz o modelo de threads e locks no cap.~8, \enquote{adapted with
minimal changes from Chapter 17 of the first edition of The Java
Language Specification}, nas seções §8.12 a §8.14. Na §8.14, um
\texttt{notify} remove do wait set uma thread escolhida arbitrariamente.

Sobre a origem, o white paper diz que as primitivas de
sincronização seguem o paradigma de monitor e variável de condição
\enquote{introduced twenty years ago by C.A.R. Hoare}, e que foram
modeladas no Cedar/Mesa (§7.1 do white paper,
\citetitle{paper:java-language-environment}). O artigo de Hoare
(\citetitle{paper:hoare1974-monitors}, 1974) \enquote{develops
Brinch-Hansen's concept of a monitor as a method of structuring an
operating system}, com variáveis de condição e as operações
\texttt{wait} e \texttt{signal}. Em 1999, Per Brinch Hansen
escreveu em \citetitle{paper:brinch-hansen1999-insecure}:
\enquote{The claim that Java supports monitors is shown to be false.}
Para ele, o que o Java oferece \enquote{is just a programming style that
imitates insecure monitors}, e o artigo conclui que o Java ignora os
últimos vinte e cinco anos de pesquisa em linguagens de programação
paralela.

\paragraph{O que mudou depois.}
Na JLS do Java SE~25, o cap.~17 continua se chamando
\enquote{Threads and Locks}, mas a sincronização está na §17.1, o wait
set na §17.2 e o modelo de memória na §17.4. A §17.1 define:
\enquote{Each object in Java is associated with a monitor, which a
thread can lock or unlock.} Só uma thread por vez detém o lock de um
monitor. A §17.1 descreve o bloco e o método \texttt{synchronized} com o texto
da §17.13 da 2ª edição, exceto que o que se trava é o monitor do
objeto e sai a observação sobre a memória principal. A mesma seção cita as classes de
\texttt{java.util.concurrent} entre os outros meios de sincronização.

A §17.2.1 detalha o \texttt{wait}. Se a thread não detém o lock do
objeto, a chamada lança \texttt{IllegalMonitorStateException}. Se
detém, ela entra no wait set e executa uma liberação para cada
travamento pendente; ao sair, refaz os mesmos travamentos. A lista de saídas do wait set cresceu: além
de \texttt{notify}, \texttt{notifyAll} e tempo limite, entram a
interrupção da thread e uma ação interna da implementação. Essa última é
o \termo{spurious wakeup}{Acordar espúrio: a thread sai de
\texttt{wait} sem \texttt{notify}, interrupção nem fim do tempo
limite.}: \enquote{Implementations are permitted, although not
encouraged, to perform \enquote{spurious wake-ups}}. A mesma seção diz que
essa permissão exige usar \texttt{wait} só dentro de laços que só
terminam quando a condição esperada vale. A §17.2.2 não dá garantia sobre qual
thread o \texttt{notify} escolhe, e, após um \texttt{notifyAll}, só uma
thread por vez retoma o lock. O javadoc de \texttt{Object.wait} no
JDK~25 (\citetitle{doc:javase25-object}) recomenda testar a condição
num laço \texttt{while} em volta da chamada.

Na JVMS do Java SE~25, a §2.11.10 diz que a JVM sincroniza métodos e
sequências de instruções com \enquote{a single synchronization
construct: the monitor}. A sincronização é feita pela entrada e saída
do monitor, explícita, com \texttt{monitorenter} e \texttt{monitorexit},
ou implícita, nas instruções de chamada e retorno de método (§3.14).

\begin{center}
\begin{tabular}{@{}p{0.22\textwidth}p{0.36\textwidth}p{0.36\textwidth}@{}}
\toprule
Aspecto & JLS/JVMS 2ª ed. (§) & SE 25 (§) \\
\midrule
Onde está & JLS §17.12--§17.14; JVMS §8.12--§8.14 & JLS §17.1 e
§17.2; JVMS §2.11.10 e §3.14 \\
Unidade & Um lock por objeto (JLS §17.13) & Um monitor por objeto
(JLS §17.1) \\
\texttt{wait} sem o lock & Só deve ser chamado com o lock (JLS §17.14)
& Lança \texttt{IllegalMonitor\-StateException} (JLS §17.2.1) \\
Saídas do wait set & \texttt{notify}, \texttt{notifyAll}, tempo limite
(JLS §17.14) & Também interrupção e ação interna (JLS §17.2.1) \\
Spurious wakeup & Não mencionado (JLS cap.~17) & Permitido, não
encorajado (JLS §17.2.1) \\
\bottomrule
\end{tabular}
\end{center}

\paragraph{Onde fica o \texttt{wait} dentro do monitor?}
O exemplo abaixo é um buffer de duas posições. \texttt{colocar} espera
enquanto o buffer está cheio, \texttt{retirar} espera enquanto está
vazio, e os dois chamam \texttt{notifyAll} depois de mudar o estado.

\lstinputlisting[language=Java]{generated/source/EsperaEmLaco.java}

\inputbytecode{EsperaEmLaco}

\begin{itemize}
\item O cabeçalho mostra \texttt{synchronized void colocar(int)}, mas o
  código não tem \texttt{monitorenter} nem \texttt{monitorexit}. O
  método leva a flag \texttt{ACC\_SYNCHRONIZED}, conferida pelas
  instruções de chamada: a thread entra no monitor de \texttt{this} antes
  do offset~0 e sai dele no retorno, normal ou por exceção (JVMS SE~25,
  §2.11.10). O bloco \texttt{synchronized}, com as duas instruções, está
  na seção~\ref{sec:sincronizacao}.
\item Em \texttt{colocar}, os offsets 0--9 testam
  \texttt{quantidade == itens.length}; o \texttt{if\_icmpne} do
  offset~9 sai do laço para o offset~19. Nos offsets 12--13,
  \texttt{aload\_0} e \texttt{invokevirtual} chamam
  \texttt{Object.wait} sobre \texttt{this}, o mesmo objeto cujo monitor
  o método detém. O \texttt{goto} do offset~16 volta ao offset~0, e a
  condição é testada de novo.
\item Dentro da chamada do offset~13, a thread executa as liberações
  e, ao sair do wait set, os travamentos de volta, antes de
  \texttt{wait} retornar (JLS SE~25, §17.2.1). Entre as liberações e
  os travamentos de volta, a thread não detém o monitor de
  \texttt{this}: o javadoc de \texttt{Object.wait}
  (\citetitle{doc:javase25-object}) diz que ela passa a
  \enquote{relinquish any and all synchronization claims on this
  object}.
\item Em \texttt{retirar}, o padrão se repete: \texttt{ifne} no
  offset~4, \texttt{wait} no offset~8 e \texttt{goto} para o offset~0
  no offset~11. \texttt{notifyAll} está no offset~51 dos dois métodos,
  antes do \texttt{return} (54) e do \texttt{ireturn} (55).
\end{itemize}

Experimento (JDK~25, imagem \texttt{eclipse-temurin:25-jdk}): a saída
real de \texttt{main} foi \texttt{1}, \texttt{2}, \texttt{3},
\texttt{4} e \texttt{5}, um número por linha.

\paragraph{O que isso muda no seu código.}
Cada recomendação tem seu escopo.
\begin{itemize}
\item \emph{JDK~25:} chame \texttt{wait} dentro de um \texttt{while}
  que testa a condição, nunca de um \texttt{if}. A JLS SE~25 permite o
  spurious wakeup (§17.2.1), e o javadoc de \texttt{Object.wait} do
  JDK~25 recomenda o laço. O javadoc já descreve o spurious wakeup e
  pede o laço desde o Java~5, de 2004 (\citetitle{doc:javase5-object}).
\item \emph{JLS 2ª ed. e JDK~25:} todo objeto tem um lock (JLS 2ª
  ed., §17.13) ou monitor (JLS SE~25, §17.1), que qualquer thread com
  uma referência a ele pode travar. O exemplo acima trava \texttt{this}
  de propósito, para mostrar um método \texttt{synchronized} sem
  \texttt{monitorenter}; a flag \texttt{ACC\_SYNCHRONIZED} só aparece
  com \texttt{javap -v}. Por que preferir um objeto de lock privado está
  na seção~\ref{sec:sincronizacao}.
\item \emph{Java~5 (2004) em diante:} para código novo, prefira as
  classes de \texttt{java.util.concurrent}, que a própria §17.1 da JLS
  SE~25 cita como alternativa de sincronização. O pacote veio da
  \citetitle{doc:jsr166}, com versão final em 30 de setembro de 2004.
\end{itemize}

\paragraph{Onde aprofundar.}
O \texttt{monitorenter} e o \texttt{monitorexit} estão na
seção~\ref{sec:sincronizacao}. As garantias de visibilidade de
\texttt{synchronized} são tema da seção~\ref{sec:jmm-historia}, e as
alternativas de \texttt{java.util.concurrent}, da
seção~\ref{sec:juc}.

\section{De threads de plataforma a threads virtuais}
\label{sec:threads-ate-vt}

\section{O modelo de memória ao longo do tempo}
\label{sec:jmm-historia}

\section{\texttt{java.util.concurrent}}
\label{sec:juc}

\section{Muitos requests, memória compartilhada, vários núcleos}
\label{sec:muitos-requests}

\paragraph{O problema.}
Num servidor de servlets, várias threads executam o mesmo código. A
\citetitle{doc:jakarta-servlet61}, §2.3.3.1, diz que
\enquote{A servlet container may send concurrent requests through the
service method of the servlet}, e o desenvolvedor do servlet precisa
prever a execução de várias threads nesse método. A mesma seção
recomenda não sincronizar o método \texttt{service}, por causa da
perda de desempenho: \enquote{It is strongly recommended that
Developers not synchronize the service method}.

O texto \citetitle{artigo:kegel-c10k}, de Dan Kegel, traz copyright de
1999 a 2018 e última atualização em 5 de fevereiro de 2014. A abertura pergunta: \enquote{It's
time for web servers to handle ten thousand clients simultaneously,
don't you think?} O texto defende que atender dez mil clientes ao mesmo
tempo pede algo além de uma thread por cliente, e diz sobre o Java:
\enquote{Up through JDK 1.3, Java's standard networking libraries
mostly offered the one-thread-per-client model.}

\paragraph{A decisão.}
A \citetitle{doc:jsr51}, com Final Release em 9 de maio de 2002, entrou
no Java~1.4 (2002). O líder da JSR foi Mark Reinhold, e Kegel e Doug
Lea estavam no grupo de especialistas. A descrição da JSR lista
\enquote{APIs for scalable I/O, fast buffered binary and character I/O,
regular expressions, charset conversion, and an improved filesystem
interface}. Segundo Kegel, o \texttt{java.nio} do JDK~1.4 trouxe ao
Java o I/O não bloqueante.

\paragraph{O que mudou depois.}
Depois veio o padrão reativo. O \citetitle{doc:reactive-streams} se
define como \enquote{a standard for
asynchronous stream processing with non-blocking back pressure}. A
\termo{back pressure}{Contrapressão: o consumidor diz ao produtor quanto
consegue receber, para não ser inundado.} é parte da definição. A
versão~1.0.0 para a JVM
(\citetitle{doc:reactive-streams-jvm}) é de 29 de abril de 2015. O
JDK~9 trouxe a interface \texttt{java.util.concurrent.Flow}
(\enquote{Since: 9}), cujo javadoc diz: \enquote{These interfaces
correspond to the reactive-streams specification}
(\citetitle{doc:javase9-flow}).

O \citetitle{doc:jep444} descreve o custo desse caminho. No estilo
assíncrono, o código devolve a thread a um pool enquanto espera I/O, e
o fim da operação é avisado por um \termo{callback}{Função passada
para outra, que a chama quando um evento acontece, como o fim de uma
operação de I/O.}. A lógica de cada requisição vira uma sequência de
estágios pequenos, em geral lambdas, compostos por uma API como
\texttt{CompletableFuture} ou um framework reativo, e o programador
abre mão de construções como laços e blocos \texttt{try}/\texttt{catch}
(\enquote{They thus forsake the language's basic sequential composition
operators}). Um dos objetivos do JEP é que servidores no estilo
thread-per-request escalem com virtual threads, finais no JDK~21
(seção~\ref{sec:threads-ate-vt}).

No Quarkus, o guia
\citetitle{doc:quarkus-reactive-architecture} diz que \enquote{A few
threads can handle many concurrent I/O}. Uma requisição HTTP é
processada numa dessas poucas threads de I/O, que por isso não podem
bloquear (\enquote{you can't block the thread anymore}). Código
bloqueante é desviado para \termo{worker threads}{Threads de um pool
separado, usadas para executar código que bloqueia sem ocupar as
threads de I/O.}: \enquote{Quarkus implements a proactor pattern that
switches to worker thread when needed}. O guia
\citetitle{doc:quarkus-virtual-threads} diz que \enquote{the support of
virtual thread is implemented using the @RunOnVirtualThread
annotation}; o que a anotação faz está na
seção~\ref{sec:virtual-threads}.

\paragraph{O que isso muda no seu código.}
Cada recomendação tem seu escopo.
\begin{itemize}
\item \emph{Jakarta Servlet~6.1, cujos containers exigem Java SE~17 ou
  mais novo (§1.2):} o mesmo servlet pode atender
  requisições concorrentes pelo método \texttt{service}, então escreva
  esse código prevendo várias threads (seções~\ref{sec:monitores}
  e~\ref{sec:jmm-historia}). A especificação recomenda não
  sincronizar o método \texttt{service}.
\item \emph{JDK~9 em diante:} \texttt{java.util.concurrent.Flow} é a
  versão do JDK das interfaces do Reactive Streams.
\item \emph{JDK~21 em diante:} código bloqueante no estilo
  thread-per-request pode escalar com virtual threads, sem reescrever
  a lógica em estágios assíncronos (seção~\ref{sec:threads-ate-vt}).
\item \emph{Quarkus~3, sobre JDK~17 ou mais novo
  (\citetitle{doc:quarkus-getting-started}):} código nas threads de I/O não pode bloquear, e
  código bloqueante vai para as worker threads; com Java~21 ou mais
  novo, \texttt{@RunOnVirtualThread} é a terceira opção
  (seção~\ref{sec:vt-vs-reativo}).
\end{itemize}

\paragraph{Onde aprofundar.}
Quando escolher virtual threads, reativo ou worker threads no Quarkus
é o tema da seção~\ref{sec:vt-vs-reativo}. As regras de uso das
virtual threads estão na seção~\ref{sec:virtual-threads}.

\section{Onde aprofundar}

As fontes principais deste capítulo, na ordem em que aparecem. Acesso
(aberto ou pago) e links também na bibliografia.

\subsection*{Especificações}
\begin{itemize}
\obra{livro:gosling2000-jls2}{A JLS 2ª ed., de onde vêm as regras de
  threads, locks e memória antes da JSR~133 (cap.~17).}
\obra{livro:lindholm1999}{A JVMS 2ª ed., com os quatro passes do
  verificador (seção~\ref{sec:verificacao}).}
\obra{doc:jls25}{A JLS atual, para comparar o cap.~17 com o da 2ª ed.}
\obra{doc:jvms25}{A JVMS atual, com o verificador por type checking e as
  flags de campo e de método mostradas nas listagens.}
\end{itemize}

\subsection*{Decisões e história}
\begin{itemize}
\obra*{paper:java-language-environment}{O white paper que lista os
  objetivos de projeto (seção~\ref{sec:objetivos-1995}).}
\obra*{paper:hoare1974-monitors}{O artigo de monitores; o white paper
  diz que o paradigma de monitor do Java vem de Hoare.}
\obra*{paper:brinch-hansen1999-insecure}{A crítica de 1999 aos monitores
  do Java (seção~\ref{sec:monitores}).}
\obra*{paper:manson-tese-jmm-2004}{A tese, aberta, que explica o que
  estava errado no modelo de memória antigo e o modelo novo.}
\obra*{doc:jsr133-faq}{Os sintomas do modelo antigo e o
  \texttt{volatile} e o \texttt{final} do novo, em perguntas e respostas.}
\obra*{paper:lea2000-forkjoin}{O framework Fork/Join e o work-stealing,
  antes de entrarem no JDK (seção~\ref{sec:juc}).}
\obra*{artigo:kegel-c10k}{O texto que pergunta como atender dez mil
  clientes ao mesmo tempo (seção~\ref{sec:muitos-requests}).}
\obra*{doc:jep444}{Por que as virtual threads existem: dar escala, não
  velocidade (seção~\ref{sec:threads-ate-vt}).}
\end{itemize}

\subsection*{Para o código do dia a dia}
\begin{itemize}
\obra*{livro:goetz2006}{Visibilidade, publicação segura e
  \texttt{java.util.concurrent} na prática.}
\obra{doc:javase25-juc}{O javadoc do pacote, com as coleções
  concorrentes e quando preferi-las a um mapa sincronizado.}
\end{itemize}

\section{Desafios práticos}

\desafio{\emph{Muitos requests.} Escreva um programa que envia 10.000
tarefas com \texttt{Thread.sleep} de um segundo para um
\texttt{Executors.newFixedThreadPool(200)} e depois para um
\texttt{Executors.newVirtualThreadPerTaskExecutor()}, no formato do
exemplo do JEP~444 (seção~\ref{sec:threads-ate-vt}). Meça o tempo total
de cada um. Troque o \texttt{sleep} por um cálculo de um segundo e meça
de novo. Anote o JDK, o número de CPUs e a máquina junto dos números.}

\desafio{\emph{Memória compartilhada: \texttt{wait} e \texttt{notifyAll}.}
Copie o \texttt{EsperaEmLaco} da seção~\ref{sec:monitores} e acrescente
um segundo consumidor. Depois troque os dois \texttt{while} por
\texttt{if}. Rode várias vezes cada versão. Usando a JLS SE~25,
§17.2.1, explique o que pode acontecer com a versão com \texttt{if} que
o laço impede.}

\desafio{\emph{Memória compartilhada: publicação.} Copie o
\texttt{PublicacaoSegura} da seção~\ref{sec:jmm-historia}, tire o
\texttt{volatile} do campo \texttt{instancia} e gere o \texttt{javap -v
-p}. Compare as flags do campo com as da listagem do capítulo. Com as
regras de \emph{happens-before} da JLS SE~25, §17.4.5, diga que garantia
o \texttt{obter()} perde e por que o \texttt{obterPorPortador()} não
depende dessa flag.}

\desafio{\emph{Vários núcleos.} Numa cópia do repositório, mude o laço
\texttt{for t in 1 8} de \texttt{scripts/bench.sh} para incluir 2, 4 e o
número de CPUs da sua máquina, e rode o benchmark \texttt{Contencao} da
seção~\ref{sec:juc}. Faça um gráfico de ops/\textmu{}s por número de threads
para os três métodos. Como cada curva muda com o número de threads? Compare com
o cabeçalho do seu \texttt{contencao.txt} (JDK, CPUs, máquina) antes de
comparar com a tabela do capítulo.}

